\section{Explicit constants and the remaining question}\label{sec:certification}

The analytic estimates were stated in terms of $\pi$, exponentials and
rational operations. Their validity does not depend on an oscillatory
quadrature or a finite scan of the spectral axis. For completeness we
describe an elementary enclosure procedure for every decimal inequality
used in the text, and give outward bounds with ample rounding margin.

\subsection{A rational enclosure procedure}

Machin's identity gives
\begin{equation}\label{eq:machin}
 \pi=16\arctan(1/5)-4\arctan(1/239).
\end{equation}
For verification, if $\alpha=\arctan(1/5)$, the double-angle formula gives
$\tan(2\alpha)=5/12$ and $\tan(4\alpha)=120/119$.
With $\beta=\arctan(1/239)$ the subtraction formula then gives
$\tan(4\alpha-\beta)=1$. The angle lies in $(0,\pi/2)$, so it is
$\pi/4$, proving \eqref{eq:machin}. For either $q=5$ or $q=239$, put
\[
 U_q=\sum_{j=0}^{40}\frac{(-1)^j}{(2j+1)q^{2j+1}},\qquad
 L_q=U_q-\frac1{83q^{83}}.
\]
The alternating-series estimate gives
$L_q\le\arctan(1/q)\le U_q$.
All four endpoints are rational. Substitution with the appropriate
signs into \eqref{eq:machin} encloses $\pi$.

For the exponential, first suppose $0\le z\le1/8$ and set
$P_{32}(z)=\sum_{j=0}^{32}z^j/j!$. The omitted positive series has first
term $z^{33}/33!$, and the ratio of each following term to its predecessor
is at most $z/34$. Therefore
\begin{equation}\label{eq:exp-enclosure}
 P_{32}(z)\le e^z\le
 P_{32}(z)+\frac{z^{33}}{33!(1-z/34)}.
\end{equation}
For a nonnegative rational interval $[a,b]$ enclosing an argument, choose
$k$ so that $b/2^k\le1/8$. Evaluate the lower polynomial at $a/2^k$ and
the upper expression at $b/2^k$, then square the positive interval $k$
times. This encloses $e^{[a,b]}$ by monotonicity. Reciprocal endpoints,
in reverse order, enclose $e^{-[a,b]}$.

Ordinary interval arithmetic now encloses the constants in
Section~\ref{sec:positivity} and the explicit tail constants of
Section~\ref{sec:shifted}. Products use the minimum and maximum of the
four endpoint products; division uses the reciprocal interval after
checking that the denominator excludes zero. All denominators used below
are strictly positive.

One may perform the calculation with exact fractions throughout.
For a bounded-size implementation, round each lower endpoint down and
each upper endpoint up to a multiple of $2^{-256}$ after every interval
operation. These outward roundings preserve containment by induction.
No machine floating-point approximation to $\pi$ or to an exponential
is needed. For the sample shift $d=\pi/16$, use
$\cos(4d)=1/\sqrt2>7/10$; the last inequality follows by squaring
positive quantities. Substitution of $7/10$ in the tail majorant is
conservative and avoids a square-root enclosure.

\subsection{Certified values}

Let $d_5=29/200$ and let $\rho_{s,5}$ denote the strip-relative constant
from Section~\ref{sec:positivity}. In the shifted-source bound set
$m=5$, $d=\pi/16$, $T=1/2$, and $Y=1$. Write $\overline K_j$ for its
conservative expression with $\cos(4d)$ replaced by $7/10$.
Thus $K_j\le\overline K_j$; the intervals in the last three rows
enclose the conservative expressions, not the exact-cosine constants.

\begin{table}[htbp]
\centering
\caption{Outward rational enclosures for the current estimates.}
\label{tab:constants}
\begin{tabular}{@{}lrr@{}}
\toprule
Quantity & Lower bound & Upper bound\\
\midrule
$\kappa$ & $0.00084868730924288$ & $0.00084868730924289$\\
$\rho_3$ & $0.00026291407856$ & $0.00026291407858$\\
$\rho_5$ & $7.25125221685\,10^{-25}$ & $7.25125221687\,10^{-25}$\\
$d_5\rho_5$ & $1.05143157144\,10^{-25}$ & $1.05143157145\,10^{-25}$\\
$d_5^2\rho_5$ & $1.52457577859\,10^{-26}$ & $1.52457577860\,10^{-26}$\\
$\rho_{s,5}$ & $8.47360190576\,10^{-25}$ & $8.47360190578\,10^{-25}$\\
$\overline K_0$ & $7.95799323606\,10^{-23}$ & $7.95799323608\,10^{-23}$\\
$\overline K_1$ & $1.60045950329\,10^{-23}$ & $1.60045950331\,10^{-23}$\\
$\overline K_2$ & $3.22054543607\,10^{-24}$ & $3.22054543609\,10^{-24}$\\
\bottomrule
\end{tabular}
\end{table}

These enclosures prove the strict rounded bounds
\[
 \rho_3<0.000263,\quad \rho_5<7.3\,10^{-25},\quad
 d_5\rho_5<1.06\,10^{-25},\quad d_5^2\rho_5<1.54\,10^{-26},
\]
\[
 \rho_{s,5}<8.5\,10^{-25},\qquad
 (\overline K_0,\overline K_1,\overline K_2)
 <(8\,10^{-23},1.61\,10^{-23},3.23\,10^{-24}),
\]
where the last inequality is componentwise.
The carrier slope bound $1/8$ and curvature budget $1/64$ follow
analytically from the disk-map argument, not from decimal certification.

The accompanying standard-library program
\texttt{certify\_updated\_manuscript.py} implements these interval rules,
replays the original manuscript constants, and verifies every endpoint
in Table~\ref{tab:constants}. Its record retains rational endpoints,
arithmetic parameters and source identities. The earlier constant
programs remain available as research provenance.
These calculations verify explicit constants conditional on their
analytic derivation. They certify neither oscillatory quadrature of the
retained shifted modes nor a positive physical-scale margin.

\subsection{The remaining mathematical question}

The finite Euler calculation yields mixed heat weights, commuting complex
shifts, and complete half-line packets with their boundary integrals.
Positive-real-part geometry now gives a zero-free arithmetic carrier
on a large strip and bounds its contribution to the normalized Laguerre
identity by $E/64$. The derivative and complex-strip remainder bounds
make precise analytic comparisons possible without an auxiliary pole.

The first mode nevertheless cancels every algebraic endpoint contribution
of the positive higher modes. At a hypothetical collision,
$H=O(s)$ and $E=O(s^2)$ along the stationary predecessor.
The scoped controls retain the stronger packet geometry, so that geometry
alone cannot supply the missing signed estimate.

Section~\ref{sec:shifted} uses complete reflection before truncation.
It supplies an explicit derivative error and an asymptotic comparison
with a nonzero envelope. Neither is a lower bound for the complete
oscillatory field. A finite exclusion needs a strict jet or contour margin
with every evaluation and whole-domain error included. A global Newman
bound additionally needs complete finite coverage and an effective outer
region, or every hypothesis of a published barrier theorem. Extracting
such a margin from the complete modular and arithmetic data remains the
substantive research problem.
